\documentclass[12pt,a4paper]{article}
\usepackage{fontspec}
\usepackage{polyglossia}
\setmainlanguage{russian}
\setotherlanguage{english}
\setmainfont{Liberation Serif}
\usepackage[left=1.3cm,right=1.3cm,top=1.0cm,bottom=0.9cm]{geometry}
\usepackage{amsmath}
\setlength{\parindent}{0pt}
\pagestyle{empty}

\newcounter{zad}

\newcommand{\shapka}[1]{%
  \noindent{\Large\bfseries Обратный ход}%
  \hfill{\small Спецмат, 7\,И \quad·\quad 15 сентября \quad·\quad вариант #1}%
  \vspace{4pt}\hrule\vspace{7pt}
  \noindent Фамилия, имя:\ \hrulefill\rule[-4mm]{0pt}{10mm}
  \vspace{2pt}\hrule
  \setcounter{zad}{0}}

% правый блок: большая рамка под ответ + узкая под подпись
\newcommand{\pole}[1]{%
  \parbox[t]{4.7cm}{%
    \fbox{\parbox[t][#1mm][t]{4.05cm}{\strut}}\\[1mm]
    \fbox{\parbox[t][4mm][t]{4.05cm}{\strut}}}}

\newcommand{\zad}[2]{%
  \stepcounter{zad}%
  \vspace{3pt}\par\noindent
  \begin{tabular}{@{}p{12.5cm}@{\hspace{5mm}}p{4.8cm}@{}}
  \textbf{\thezad.}\ #2 & \pole{#1}
  \end{tabular}\par}

\newcommand{\zadz}[3]{%
  \vspace{3pt}\par\noindent
  \begin{tabular}{@{}p{12.5cm}@{\hspace{5mm}}p{4.8cm}@{}}
  \textbf{#1$^{*}$.}\ #3 & \pole{#2}
  \end{tabular}\par}

\begin{document}

%======================= ВАРИАНТ A =======================
\shapka{A}

\zad{11}{Кира задумала число, прибавила к нему~5, разделила сумму на~3, умножила результат на~4, вычла~6, разделила на~7 и получила~2. Какое число задумала Кира?}

\zad{11}{В чашке с питательной смесью живут бактерии, и каждый час их становится вдвое больше. В 20:00 бактерий оказалось 1600. Во сколько их было 400?}

\zad{13}{Пятеро друзей ели чипсы из одной большой пачки. Каждый по очереди подходил к пачке и съедал ровно половину того, что в ней оставалось. Пятый съел 50~граммов. Сколько чипсов было в пачке сначала?}

\zad{11}{Улитка забирается на дерево высотой 17~метров. За день она поднимается на 4~метра, а за ночь сползает вниз на~3. В какой по счёту день улитка доберётся до вершины?}

\zad{14}{В компьютерной игре есть портал: когда герой проходит через него, число его монет удваивается, но страж портала сразу забирает 24~монеты. Герой прошёл через портал три раза, и монет у него не осталось. Сколько монет было у героя вначале?}

\zad{21}{Аня, Боря и Валя играли в настольную игру на фишки. Проигравший в партии отдаёт каждому из двух других столько фишек, сколько у того уже есть. Сыграли три партии, и каждый ровно один раз проиграл: сначала Аня, потом Боря, потом Валя. В конце у каждого оказалось по 24~фишки. Сколько фишек было у каждого перед началом игры?}

\zad{14}{Летела стая гусей. На первом озере села половина стаи и ещё полгуся, на втором~--- половина оставшихся и ещё полгуся, и так дальше. На седьмом озере сели все оставшиеся гуси. Сколько гусей было в стае?}

\zadz{8}{13}{На экране написано число~458. За один ход разрешается либо удвоить число на экране, либо стереть у него последнюю цифру. Как за несколько ходов получить число~14?}

\zadz{9}{18}{Преподаватель дал на уроке несколько задач. За первые пятнадцать минут школьники решили две трети всех задач и ещё две трети задачи. За вторые пятнадцать минут~--- три четверти оставшихся задач и ещё три четверти задачи. Под конец они дорешали всё: половину оставшихся задач и ещё половину задачи. Сколько задач дал преподаватель?}

\zadz{10}{14}{У Миши калькулятор с двумя кнопками: одна удваивает число на экране, другая переставляет его цифры как угодно (ноль на первом месте оказаться не может). Сейчас на экране~1. Сможет ли Миша получить 74? А 112?}

\newpage

%======================= ВАРИАНТ B =======================
\shapka{B}

\zad{11}{Кирилл задумал число, вычел из него~4, умножил результат на~3, прибавил~9, разделил на~6, умножил на~5 и получил~20. Какое число задумал Кирилл?}

\zad{11}{На пруду растут кувшинки, и каждую ночь их становится вдвое больше. К 15~июля кувшинок стало 800. Какого числа их было 100?}

\zad{13}{Пятеро друзей ели орешки из одной большой пачки. Каждый по очереди подходил к пачке и съедал ровно половину того, что в ней оставалось. Пятый съел 30~граммов. Сколько орешков было в пачке сначала?}

\zad{11}{Робот-мойщик забирается на мачту высотой 23~метра. За день он поднимается на 6~метров, а за ночь соскальзывает вниз на~5. В какой по счёту день робот доберётся до вершины?}

\zad{14}{В компьютерной игре есть портал: когда герой проходит через него, число его монет удваивается, но страж портала сразу забирает 32~монеты. Герой прошёл через портал три раза, и монет у него не осталось. Сколько монет было у героя вначале?}

\zad{21}{Аня, Боря и Валя играли в настольную игру на фишки. Проигравший в партии отдаёт каждому из двух других столько фишек, сколько у того уже есть. Сыграли три партии, и каждый ровно один раз проиграл: сначала Аня, потом Боря, потом Валя. В конце у каждого оказалось по 32~фишки. Сколько фишек было у каждого перед началом игры?}

\zad{14}{Летела стая гусей. На первом озере села половина стаи и ещё полгуся, на втором~--- половина оставшихся и ещё полгуся, и так дальше. На шестом озере сели все оставшиеся гуси. Сколько гусей было в стае?}

\zadz{8}{13}{На экране написано число~926. За один ход разрешается либо удвоить число на экране, либо стереть у него последнюю цифру. Как за несколько ходов получить число~14?}

\zadz{9}{18}{Преподаватель дал на уроке несколько задач. За первые пятнадцать минут школьники решили две трети всех задач и ещё две трети задачи. За вторые пятнадцать минут~--- три четверти оставшихся задач и ещё три четверти задачи. Под конец они дорешали всё: две трети оставшихся задач и ещё две трети задачи. Сколько задач дал преподаватель?}

\zadz{10}{14}{У Миши калькулятор с двумя кнопками: одна удваивает число на экране, другая переставляет его цифры как угодно (ноль на первом месте оказаться не может). Сейчас на экране~1. Сможет ли Миша получить 38? А 228?}

\end{document}
